% Note: the official solution is divided into Part (a) and Part (b) only;
% its Part (b) covers both parts (b) and (c) of the problem statement.
\begin{description}
\item[Part (a)] To determine whether the system can be detected by LISA, we
  need to determine two quantities: the gravitational-wave frequency and the
  characteristic strain amplitude.

  It is straightforward to calculate the gravitational-wave frequency:
  \[ f = \frac{2}{P} = \frac{2}{414.79} = 4.8 \times 10^{-3}\,\mathrm{Hz}. \]
  This frequency is within the LISA band and close to the maximum LISA
  sensitivity. \hfill(2 points)

  To estimate the characteristic strain amplitude, we need to know the chirp
  mass $\mathcal{M}$. The other required quantities (such as the gravitational
  wave frequency, distance, and duration of the LISA observations) are already
  known.

  The orbital period change rate $\Delta P/\Delta t$ is given in the problem.
  We need to link it to $\Delta a/\Delta t$ which we know from Note 3 is
  linked to the mass function. From Kepler's third law, we know that:
  \[ \frac{a^3}{P^2} = \frac{G(M_1 + M_2)}{4\pi^2}, \]
  where $M_1$ and $M_2$ are masses of both components of the system. If, due
  to the emission of gravitational waves, the semi-major axis changes from $a$
  to $a + \Delta a$, then the orbital period changes from $P$ to
  $P + \Delta P$, as the masses of both components are constant. Then:
  \[ \frac{a^3}{P^2} = \frac{(a+\Delta a)^3}{(P+\Delta P)^2}
     = \frac{a^3(1+\Delta a/a)^3}{P^2(1+\Delta P/P)^2}
     = \frac{a^3}{P^2}\left(1 + \frac{3\Delta a}{a}
       - \frac{2\Delta P}{P}\right), \]
  hence:
  \[ \frac{\Delta P}{P} = \frac{3}{2}\,\frac{\Delta a}{a}. \]
  Here, we used the fact that $(1+x)^n \approx 1 + nx$ for $x \ll 1$.

  Therefore:
  \begin{align*}
    \frac{\Delta P}{\Delta t}
      &= \frac{3P}{2a}\frac{\Delta a}{\Delta t}
       = -\frac{3}{2}\cdot\frac{64}{5}\,\frac{P}{a}\,\frac{G^3}{c^5}\,
         \frac{M_1 M_2 (M_1+M_2)}{a^3} \\
      &= -\frac{96}{5}\,\frac{G^3}{c^5}\,\frac{P}{a}\,
         \frac{M_1 M_2 (M_1+M_2)}{G(M_1+M_2)P^2}\cdot 4\pi^2 \\
      &= -\frac{96}{5}(2\pi)^2\,\frac{G^2}{c^5}\,\frac{M_1 M_2}{aP}
       = -\frac{96}{5}(2\pi)^2\,\frac{G^2}{c^5}\,\frac{M_1 M_2}{P}\,
         \frac{(2\pi)^{2/3}}{G^{1/3}(M_1+M_2)^{1/3}P^{2/3}} \\
      &= -\frac{96}{5}(2\pi)^{8/3}\,\frac{G^{5/3}}{c^5}\,
         \frac{M_1 M_2}{(M_1+M_2)^{1/3}}\,\frac{1}{P^{5/3}} \\
      % note: the source prints the expression on the line above twice.
      &= -\frac{96}{5c^5}(2\pi)^{8/3}
         \left(\frac{G\mathcal{M}}{P}\right)^{5/3}
       = -\frac{192\pi}{5c^5}(G\mathcal{M})^{5/3}
         \left(\frac{P}{2\pi}\right)^{-5/3},
  \end{align*}
  Thus, by knowing the orbital period and its rate of change from
  observations, we can determine the chirp mass:
  \[ \mathcal{M} = \left(\frac{5}{192\pi}\right)^{3/5}\frac{c^3}{G}\,
     \frac{P}{2\pi}\left(-\frac{\Delta P}{\Delta t}\right)^{3/5}
     = 0.319\,M_\odot, \]
  and find the characteristic strain amplitude $h$.

  Alternatively, if we notice that the characteristic strain amplitude $h$
  depends on $(G\mathcal{M})^{5/3}$ which we can get directly from the rate of
  change of the period:
  \[ (G\mathcal{M})^{5/3} = -\frac{\Delta P}{\Delta t}
     \left(\frac{P}{2\pi}\right)^{5/3}\frac{5c^5}{192\pi}, \]
  we can skip calculating the chirp mass and get the characteristic strain
  amplitude directly, which is all we need to check if the binary can be
  detected.

  Either way, the gravitational wave strain is:
  \[ h = -\frac{5c}{192\pi^2}\,P\,\frac{\Delta P}{\Delta t}\,\frac{1}{D}. \]
  If we plug in the numerical values, we get $h = 1.1 \times 10^{-22}$ and
  $S = 8.4 \times 10^{-20}$. Checking the plot, this is above the expected
  sensitivity of LISA at 5\,mHz. Thus, this object should be detected by LISA.

\item[Part (b)] To determine the masses of both components $M_1$ and $M_2$ we
  need two simultaneous $M_1 - M_2$ relations. The expression for the chirp
  mass:
  \[ \mathcal{M} = \frac{(M_1 M_2)^{3/5}}{(M_1 + M_2)^{1/5}}, \]
  gives us one relation, and we can calculate the chirp mass of the observed
  system from the rate of change of the period, if that was not already done
  in part (a):
  \[ \mathcal{M} = \left(\frac{5}{192\pi}\right)^{3/5}\frac{c^3}{G}\,
     \frac{P}{2\pi}\left(-\frac{\Delta P}{\Delta t}\right)^{3/5}
     = 0.319\,M_\odot. \]
  The second relation can be obtained from the mass--radius relation for white
  dwarfs given in the table and Kepler's third law.

  We are told that for one of the components (call it `1'),
  \[ a = \frac{R_1}{0.139}. \]
  From Kepler's third law we know that:
  \[ M_1 + M_2 = \left(\frac{a}{1\,\mathrm{au}}\right)^3
     \left(\frac{P}{1\,\mathrm{yr}}\right)^2, \]
  which will give us the mass of the second component $M_2$ from the mass of
  the first, $M_1$.

  From here, it is not possible to derive an analytical formula for the masses
  of the components. Instead, we need to use numerical methods to estimate the
  result.

  Taking the masses and radii listed in the given table as $M_1$ and $R_1$, we
  can obtain $a$ and thus $M_1 + M_2$, $M_2$ and finally $\mathcal{M}$ for
  each mass:

  \begin{center}
  \begin{tabular}{cccccc}
    \hline
    $M_1$ ($M_\odot$) & $R$ ($R_\odot$) & $a$ ($R_\odot$) & $M_1+M_2$ &
    $M_2$ & $\mathcal{M}$ \\
    \hline
    0.48 & 0.0144 & 0.104 & 0.647 & 0.167 & 0.240 \\
    0.50 & 0.0147 & 0.106 & 0.689 & 0.189 & 0.261 \\
    0.52 & 0.0150 & 0.108 & 0.732 & 0.212 & 0.283 \\
    0.54 & 0.0153 & 0.110 & 0.776 & 0.236 & 0.306 \\
    0.56 & 0.0156 & 0.112 & 0.823 & 0.263 & 0.329 \\
    0.58 & 0.0159 & 0.114 & 0.871 & 0.291 & 0.354 \\
    0.60 & 0.0162 & 0.117 & 0.922 & 0.322 & 0.379 \\
    0.62 & 0.0165 & 0.119 & 0.974 & 0.354 & 0.405 \\
    0.64 & 0.0168 & 0.121 & 1.028 & 0.388 & 0.431 \\
    \hline
  \end{tabular}
  \end{center}

  The actual chirp mass is $\mathcal{M} = 0.319\,M_\odot$. Therefore, by
  linear interpolation or graphically, we estimate $M_1 = 0.55\,M_\odot$ and
  $M_2 = 0.25\,M_\odot$.

  (The student does not need to calculate $\mathcal{M}$ for all values of
  $M_1$.)

  % FIGURE NEEDED: graphical solution -- plot of chirp mass \mathcal{M}
  % (0.24-0.43 M_sun) versus M_1 (0.48-0.64 M_sun), nine points on a nearly
  % straight rising curve; official 2023 solutions, p. 39.
\end{description}
